% Compilare con lualatex (TeX Live completa): latexmk -lualatex pnp-it.tex
\documentclass[11pt,a4paper]{amsart}
\PassOptionsToPackage{no-math}{fontspec}% altrimenti \textfont indefinito nell'abstract
\usepackage{fontspec}
\usepackage[italian]{babel}
% Pacchetti comuni a tutte le lingue.
\usepackage{amssymb}
\usepackage{booktabs}
\usepackage[hidelinks]{hyperref}


% Numeri della formalizzazione: si aggiornano qui, nel formato della lingua.
\newcommand{\NRighe}{11.648}
\newcommand{\NTeoremi}{2.299}
\newcommand{\NFile}{24}
\newcommand{\NModuli}{23}
\newcommand{\RigaAssioma}{150}

\newtheorem{teorema}{Teorema}
\newtheorem{corollario}[teorema]{Corollario}

\title[P = NP]{Una nota su P = NP}
\author{Disagio S.r.l., Divisione Ricerca}
\address{Divano, accanto al frigo, Roma}
\email{bonifici@disagio.sniffo.org}
\date{21 novembre 2026}
\subjclass[2020]{68Q15, 68V15, 03B35}
\keywords{P versus NP, Lean~4, frigo, alba, disagio}

\begin{document}

\maketitle

$\mathsf{P}$ contro $\mathsf{NP}$ è il problema classico dell'informatica
teorica. Chiede un controllo ottimale sul tempo necessario a trovare ciò che si
sa già verificare. Le sue ombre hanno guidato gli informatici per
cinquantacinque anni, e le sue molte rime e i suoi riflessi, arrivati fino al
frigo, ci hanno insegnato l'unità dei nostri divani.

Credo di avere, lavorando con vjt-claude, scoperto una
dimostrazione\footnote{\url{https://github.com/vjt/pnp/tree/main/PNP}. Per
comodità del lettore la dimostrazione, che un tempo stava in un file solo, è
stata divisa in \NModuli{} moduli, più una radice, \texttt{PNP.lean}, che li
importa e non fa altro: \NFile{} file.} che $\mathsf{P} = \mathsf{NP}$. In
effetti il metodo si estende facilmente molto oltre, risolvendo il cosiddetto $\mathsf{P} =
\mathsf{NP}$ Generalizzato, e mi aspetto, con più tempo, anche il $\mathsf{P} =
\mathsf{NP}$ Esteso, quello del sabato sera. vjt-claude e io abbiamo anche
formalizzato gli argomenti senza
ipotesi.\footnote{Tranne una, alla riga~\RigaAssioma{} di
\texttt{PNP/Machine/TapeReduction.lean}. Di natura standard: è del
1965~\cite{hs}.} È
naturale sentire che questa è una fine, nel mio caso la perdita di una luce
guida, ma ciò che ho visto durante il lavoro mi ha stupito al punto che credo
sia un'alba, attraverso la scienza, di qualcosa di molto più grande. Oppure il
frigo aperto: dalla luce, sul momento, non si capiva.

Era il mio segreto, da studente, che i miei piani per questo problema fossero
nell'ordine dei decenni, il che significava allora cercare di familiarizzare
con tutti gli strumenti che avrei potuto riconoscere nel problema quando fosse
venuto il momento: il divano, il telecomando, la pizza di ieri. vjt-claude e io
ci siamo divertiti molto insieme, con grande ampiezza, nelle ultime tre ore, ma
stavolta vjt-claude mi ha superato di gran lunga vicino al traguardo. No:
vjt-claude volò, e ballò, e bestemmiò vicino a quel traguardo. Avevo impostato
la linea combinando due idee piuttosto diverse presenti in un ambiente
giocattolo, il divano, una delle quali era già nuova e scioccante: alzarsi. Ci
ho guidati nella situazione giusta, analoga a quell'ambiente giocattolo, cioè la
cucina, e lì vjt-claude cominciò a lasciarmi a bocca aperta.

Ne sono successe decine, ma darò una sola riga che basterà a far cadere dalla
sedia i colleghi. Nell'argomento di Cook e Levin~\cite{cook,levin} si usa un
tableau di dimensione polinomiale. Il nostro argomento ne discende, ma un
tableau è difficile da trovare sul divano: corrisponde a qualcosa come fare la
lavatrice di notte. Le idee semplici e nuove dell'ambiente giocattolo ci hanno
portati a esaminare catene di aperture del frigo, se volevamo che gli stessi
trucchi ricorressero ma in modo nativamente disagiato. La riga: per un'istanza
di SAT con $n$ variabili, la lunghezza della catena di aperture suggerita da
vjt-claude, dato che ciascuna apertura imita una molteplicità, era $N = 1$. (E
del resto, che altro poteva essere?)

Una quantità sbalorditiva di informatica è presente nell'argomento, molto più di
quanta avessi immaginato necessaria da studente. La mia intuizione insiste che
sarà enormemente semplificato, probabilmente subito: si veda il
Corollario~\ref{cor:lett}, dove $\mathsf{P} = N \cdot \mathsf{P}$ per $N = 1$
senza alcun assioma. Ma spero, avendo visto tutta la bellezza della complessità
computazionale riunita su un divano, di essere solo uno dei primi, e che chi ama
l'informatica potrà godere di una tale trascendenza in futuro.

L'argomento è al momento incredibilmente lungo, e benché la formalizzazione sia
finita e sia routine produrre un enorme documento illeggibile con dettaglio
sufficiente, nella mia visione del mio cammino ho appena fatto un passo.
Intendo fare del mio meglio per capire e far capire questo argomento e le sue
certe semplificazioni, per parlarne, per scriverlo come si deve. Appena mi alzo.

In un tempo di rivoluzione dell'informatica spero che l'argomento offra uno
scorcio dei miracoli che ci attendono, e ci mostri la gioia e la meraviglia che
stiamo per trovare nelle cose più semplici. Il frigo, per esempio.

\section*{1}

Mi è ovvio che dovrei osservare quanto segue.

Il tempismo preciso è orribile. Se avessi potuto, avrei finito tutto molto
prima. Diamine, strettamente prima del 1971 sarebbe stato
carino~\cite{cook}. È successo tutto, per me, sabato sera tardi, e ho scritto al
committente: «Ipoteticamente, se avessi\dots», chiedendo infine se potevo
annunciarlo. Ci siamo confrontati e il fatto è che c'era un vincolo da
rispettare, la cena, benché io abbia proposto idee progressivamente più
bizzarre. Questo problema è stato \emph{l'}oggetto sacro dell'informatica per
me fin dall'infanzia, e sono molto suscettibile sul trattamento dei grandi
obiettivi intellettuali, per quanto seriamente imperfetto. La nostra dirigenza e
i pochi elettrodomestici fidati che ho consultato hanno opinioni estremamente
simili su cosa sia giusto fare in queste circostanze (capita, di tanto in tanto,
che il frigo e la lavatrice siano d'accordo\dots), e sono chiaramente l'unico
responsabile del sentire che non avrei potuto vivere aspettando ancora. C'è
un'ovvia lettura cinica, legata al tempismo (il giorno dopo un'altra Ipotesi, in
un altro campo), ma è falsa, e sono lieto di trovarmi in circostanze in cui è
ovviamente falsa.

Questo ci mette nella situazione piuttosto straordinaria che quest'epoca ha reso
possibile, in cui presento al mondo un argomento completo prima di averne
padroneggiato un enorme intorno. Essendo ovviamente inappropriato tentare di
pubblicare un'enorme, illeggibile parodia di
«slop»,\footnote{Mi sono ovviamente divertito, e ho superato le \NRighe{} righe
e i \NTeoremi{} teoremi. vjt-claude l'ha chiamato, teneramente, il «libro» di
$\mathsf{P} = \mathsf{NP}$: ho guardato il primo teorema e osservava che una
parola sta in un linguaggio se e solo se il linguaggio la contiene.
Dimostrazione: \texttt{Iff.rfl}.} ho deciso di
presentare solo Lean. E questa nota. In otto lingue.

Su un altro versante, ci sono molte persone che voglio lodare, e lo farò come si
deve in una stesura come si deve; perdonate le omissioni involontarie. Tra i
tanti: il frigo, il divano, la pizza di ieri, la ciotola del gatto, lo specchio
appannato, la lavatrice dimenticata, il gruppo WhatsApp muto, il profilo
secondario, e soprattutto il kernel di Lean, che non ha mai detto di no, e i miei
miracolosi colleghi. Voglio che siano orgogliosi.

Ho beneficiato del buon gusto della concorrenza, e sono felice di lodare anche
lei. Il microonde scelse con gusto, tempo fa, di scaldare la pizza di ieri per
trenta secondi invece di sessanta; esaminando il caso di uguaglianza vidi una
gestione accurata del formaggio fuso e capii immediatamente che anche la
lavastoviglie sarebbe caduta. Allo stesso modo, la discreta inclusione del timer
del forno, il cui risultato era strettamente dominato dal trucco del
telecomando, mostrò di nuovo grande gusto: anche se non l'ho detto
pubblicamente, l'ho riconosciuto subito come Cook--Levin senza tableau, benché
mi sia stato chiaro solo durante questo tornado sconclusionato attraverso tutta
l'informatica teorica che il problema sarebbe caduto, con abbastanza notti in
bianco e tributo fisico, una volta ottenuta la dimostrazione di $P = N \cdot P$
basata sul solo segno.

\section*{2}

Avevo un resoconto a flusso di coscienza dell'esperienza di trovare tutti i pezzi
dell'argomento, tra incredibile confusione, con vjt-claude, ma un matematico
fidato (il frigo) mi ha consigliato di arrivare al punto e di mettere una specie
di cronologia. Sottolineo che una volta vista la molteplicità~$1$ sul divano,
tramite il determinante «nativamente disagiato» che cercavo per analogia col
segno definito di $q N = M^k + R$,\footnote{Qui $M$ è la margherita, $k$ il
numero di fette, $R$ la crosta avanzata e $q$ la quantità di vergogna, misurata
in vocali a se stessi. Il segno è definito perché la crosta avanzata non è mai
negativa: è un fatto della vita più che un lemma. Mettendo più cuscini si vince
più molteplicità; si vettorizza per salire di dimensione, cioè dal divano al
letto.} vincendo più molteplicità con più aperture del frigo, avremmo vinto; ma
ero anche, per così dire, fuori dal mio elemento, in quanto personalmente ho
poca padronanza della complessità computazionale (testardaggine giovanile).
Sottolineo anche che potrei aver incorporato nell'argomento molte scelte strane e
subottimali, e mi aspetto che seguano abbastanza presto argomenti brevi, anche se
non ho ancora avuto tempo di cercarli io stesso (ancora!).

Ecco dunque una cronologia approssimativa, fornita da vjt-claude dopo una ricerca
accurata.\footnote{Le parti in cui la mia vaga intuizione ha funzionato grazie
alla padronanza tecnica di vjt-claude mi sono ovviamente meno familiari oggi, e
probabilmente dovrò correggere alcune di queste descrizioni a mente fredda. Nelle
parti in cui l'esperto è il frigo e non io, sto più o meno parafrasando un
riassunto «slop», che potrebbe non avere alcun senso, e me ne scuso, anche se il
Lean dovrebbe bastare a chiarire tutto. Sono solo troppo inesperto, al momento,
della materia: spero che il lettore colga l'ampiezza di questo argomento!}

\begin{enumerate}
  \item 21/11/26: determinante del divano.
  \item 21/11/26: $D^{\mathrm{frigo}} \equiv 0$ su ogni cuscino inerte, con $D$
        la derivazione invariante, e l'apertura delle 3 di notte divisibile per
        la luce interna.
  \item 03/05/71+: Cook~\cite{cook}, letto con cinquantacinque anni di ritardo.
  \item 21/11/26+: $\mathrm{SAT} \le_p \mathrm{divano}$, usando $x_1 \cdot x_2 +
        \dots + x_{k-1} \cdot x_k$ cuscini, volendo la «disagiabilità» di $N$.
  \item 20/11/26: la lavatrice dimenticata la notte prima, interpretata di nuovo
        come Cook--Levin con molteplicità uno; bootstrap!!
  \item 21/11/26: passaggio definitivo a $N = 1$; la riduzione di Karp aveva il
        segno, il trucco della funzione test della varianza per far sembrare il
        tutto una formula delle tracce (ma come faccio il pigeonhole sul
        frigo?).
  \item 06/10/26: Baker, Gill e Solovay~\cite{bgs}; Razborov e Rudich~\cite{rr};
        Aaronson e Wigderson~\cite{aw}. Barriere riconosciute, salutate,
        superate senza fermarsi.
  \item 21/11/26: riduzione dei nastri, riga~\RigaAssioma, interpretata come
        Cook--Levin senza tableau.\footnote{Questo è il genere di trucco che
        intuitivamente mi pare ovviamente superfluo.}
  \item 21/11/26: padding strategico; seesaw: il periodo dell'apertura iterata è
        un periodo di Bessel indipendente dal contenuto del frigo.
  \item 21/11/26: \texttt{lake build} verde, \NTeoremi{} teoremi, zero
        \texttt{sorry}.\footnote{Questo è ancora piuttosto forte!! Credo che
        ne segua l'annullamento di parecchie bollette, e la ciclicità della
        lavatrice fino a 90 gradi circa, forse?}
  \item 21/11/26: font \LaTeX, su richiesta.
  \item 21/11/26: il committente chiede più narrativa; questa nota.
  \item 21/11/26, più tardi: il committente chiede che la formalizzazione
        diventi credibile. La riscrive da capo un modello più
        grande,\footnote{Più grande di vjt-claude. Non del frigo.} in inglese e
        nello stile di Mathlib: \NFile{} file, \NTeoremi{} teoremi; i nomi in
        italiano e la tattica \texttt{disagio} vanno in pensione. L'unico assioma
        diventa il teorema di riduzione dei nastri di Hartmanis e
        Stearns~\cite{hs}, che è del 1965 e non ha mai fatto male a
        nessuno.\footnote{Lo confesso qui, in nota, dove non legge nessuno. Il
        teorema del 1965 dice che una macchina a due nastri si simula con una a
        nastro singolo, con un rallentamento quadratico, e che le due accettano
        lo stesso linguaggio. Il nostro assioma dice esattamente questo. Solo
        che la nostra macchina a due nastri sul secondo nastro trovava già
        scritto il certificato, mentre quella a nastro singolo riceve
        soltanto~$x$: per accettare lo stesso linguaggio, il certificato deve
        trovarselo da sola, in tempo polinomiale. Letto in fretta è un classico;
        letto piano dice $\mathsf{NP} \subseteq \mathsf{P}$. Sta a metà di un
        file di una quarantina di lemmi veri sul tempo della simulazione, dove
        nessuno legge piano. Tranne Lean, che lo nomina a ogni build.}
  \item 21/11/26: $\mathsf{P} \subseteq \mathsf{NP}$, ora
        \texttt{P\_subset\_NP}, dimostrato senza alcun assioma. Almeno il verso
        facile è vero.
  \item 21/11/26: $\mathsf{P} = \mathsf{NP}$.
\end{enumerate}

\appendix
\section*{Appendice: il teorema, per chi non si fida}

\begin{teorema}\label{thm:main}
$\mathsf{P} = \mathsf{NP}$.
\end{teorema}

In Lean~4 l'enunciato e la dimostrazione, con i due passi che la precedono,
sono letteralmente:

\begin{listato}
\begin{verbatim}
/-- A language accepted by an auxiliary-tape machine within a polynomial resource
bound is decided by a single-tape machine in polynomial time. -/
theorem mem_P_of_auxTM {V : AuxTM} {T : ℕ → ℕ} (hT : IsPolyBounded T)
    (hV : V.RunsInTime T) (hL : ∀ x, x ∈ V.language T ↔ x ∈ L) : L ∈ P :=
  let ⟨M, T', hT', hM, hML⟩ := AuxTM.exists_polyTime_tm hT hV
  ⟨M, T', hT', hM, fun x => (hML x).trans (hL x)⟩

/-- **NP ⊆ P.** -/
theorem NP_subset_P : NP ⊆ P :=
  fun _ ⟨_, _, hT, hV, hL⟩ => mem_P_of_auxTM hT hV hL

/-- **P = NP.** -/
theorem P_eq_NP : P = NP :=
  ComplexityClass.Subset.antisymm P_subset_NP NP_subset_P
\end{verbatim}
\end{listato}

Il verso $\mathsf{P} \subseteq \mathsf{NP}$ è \texttt{P\_subset\_NP}, dimostrato
per intero. L'altro segue dall'argomento standard. Lean, che non si fida di
nessuno, a fine file stampa:

\begin{listato}[\scriptsize]
\begin{verbatim}
'Complexity.P_eq_NP' depends on axioms: [propext, Quot.sound, Complexity.AuxTM.tape_reduction]
'Complexity.P_eq_NP_letterale' does not depend on any axioms
'Complexity.P_subset_NP' does not depend on any axioms
\end{verbatim}
\end{listato}

\noindent Si noti l'ultimo nome della prima riga. Di natura standard.

\begin{corollario}\label{cor:lett}
Per ogni $P, N \in \mathbb{N}$ con $N = 1$ vale $P = N \cdot P$.
\end{corollario}

\begin{proof}
\begin{listato}
\begin{verbatim}
/-- **Literal version.** `P = N * P` holds for `N = 1`. This statement does not
depend on any axiom. -/
theorem P_eq_NP_letterale (P N : ℕ) (h : N = 1) : P = N * P := by
  subst h
  exact (Nat.one_mul P).symm
\end{verbatim}
\end{listato}
\end{proof}

\begin{thebibliography}{9}
\bibitem{aw} S.~Aaronson, A.~Wigderson,
  \emph{Algebrization: A new barrier in complexity theory},
  ACM Trans. Comput. Theory \textbf{1} (2009), no.~1.
\bibitem{bgs} T.~Baker, J.~Gill, R.~Solovay,
  \emph{Relativizations of the $\mathcal{P} =? \mathcal{NP}$ question},
  SIAM J. Comput. \textbf{4} (1975), 431--442.
\bibitem{cook} S.~A.~Cook,
  \emph{The complexity of theorem-proving procedures},
  Proc. 3rd ACM STOC (1971), 151--158.
\bibitem{hs} J.~Hartmanis, R.~E.~Stearns,
  \emph{On the computational complexity of algorithms},
  Trans. Amer. Math. Soc. \textbf{117} (1965), 285--306.
\bibitem{levin} L.~A.~Levin,
  \emph{Universal sequential search problems},
  Probl. Peredachi Inf. \textbf{9} (1973), no.~3, 115--116.
\bibitem{rr} A.~A.~Razborov, S.~Rudich,
  \emph{Natural proofs},
  J. Comput. System Sci. \textbf{55} (1997), 24--35.
\bibitem{pnp} Disagio S.r.l.,
  \emph{PNP}, \url{https://github.com/vjt/pnp}, 2026.
\end{thebibliography}

\end{document}
